\section*{Chapitre \ref{ch:ll:order-gateway-and-order-management} --- Passerelle d'ordres et gestion des ordres}

\begin{solution}{exo:ll:order-gateway-and-order-management:1}
L'exécution le fait passer de \emph{en attente d'annulation} à \emph{exécuté}~; la réponse «~trop tard~» le laisse \emph{exécuté}. Avant l'exécution
arrivée, la passerelle comptait toute la quantité restante (l'annulation pouvait échouer)~; après elle, il ne reste rien.
\end{solution}

\begin{solution}{exo:ll:order-gateway-and-order-management:2}
À 0~: deux jetons, donc les deux premiers passent et le troisième est refusé. À \qty{1}{\milli\second}, le seau contient 0.3 d'un
jeton~: refusé. À \qty{4}{\milli\second}, 1.2~: il passe, en laissant 0.2. À \qty{5}{\milli\second}, 0.5~: refusé. Donc
passe, passe, refusé, refusé, passe, refusé.
\end{solution}

\begin{solution}{exo:ll:order-gateway-and-order-management:3}
L'exposition longue dans le pire cas est $200 + 100 + 100 = 400$. Le remplacement ferait passer la quantité restante la plus défavorable d'un ordre de 100 à
300, ce qui ajoute 200, soit 600 au-dessus de la limite de 500~: refusé. La plus grande quantité nouvelle acceptable ajoute 100~: un achat de 200.
\end{solution}

\begin{solution}{exo:ll:order-gateway-and-order-management:4}
$1 - e^{-50 \times 350 \times 10^{-6}} = 1 - e^{-0.0175} \approx 1.7~\%$, soit environ 347 courses sur 20\,000 annulations (si les ordres
sont restés au moins \qty{350}{\micro\second} au carnet avant d'être annulés).
\end{solution}

\begin{solution}{exo:ll:order-gateway-and-order-management:5}
Les deux exécutions ont le même numéro d'appariement~: il s'agit d'une seule exécution reçue deux fois, et la passerelle, avec 60, a raison.
La réconciliation doit indexer les exécutions par leur numéro d'appariement et compter chacune une seule fois~; un doublon mérite lui-même d'être signalé,
puisque le \omterm{def:ll:order-gateway-and-order-management:dropcopy}{drop copy} est censé constituer le registre indépendant.
\end{solution}

\begin{solution}{exo:ll:order-gateway-and-order-management:6}
Environ \qty{1.6}{\micro\second} par nouvel ordre, soit quelque trois cents fois le travail propre de la passerelle, de 4 à \qty{5}{\nano\second} par
message, et plus que le reste du parcours de l'ordre dans la machine. Des sommes cumulées le suppriment.
\end{solution}

\begin{solution}{exo:ll:order-gateway-and-order-management:7}
Chaque ordre ouvert du côté concerné entre en \emph{en attente d'annulation} à l'envoi de la demande, qui consomme un jeton du limiteur de messages (throttle)~;
l'exposition continue de compter chacun en entier jusqu'à l'arrivée de sa propre annulation ou d'une exécution. Tester contre le moteur d'appariement avec des exécutions
croisant l'annulation de masse (mass cancel)~: chaque ordre doit finir exécuté ou annulé, le \omterm{def:ll:order-gateway-and-order-management:dropcopy}{drop copy} doit se réconcilier, et la position doit rester
tout au long sous la limite.
\end{solution}

\begin{solution}{exo:ll:order-gateway-and-order-management:8}
Une exécution qui a croisé l'annulation arrive pour un ordre que la passerelle ne connaît plus~: elle est ignorée (ou traitée comme une erreur), si bien que
la position est fausse et l'exposition sous-estimée, ce qui correspond exactement à la double position de l'ouverture du chapitre. Un ordre
reste dans la table jusqu'à un rapport final (exécuté, annulé, rejeté), et la table est dimensionnée en conséquence.
\end{solution}

\begin{solution}{pb:ll:order-gateway-and-order-management:1}
\begin{enumerate}
\item Une exécution avant $c - d$ a été signalée avant $c$, si bien que la stratégie la connaissait au moment d'annuler~; une exécution après $c + L$ rencontre un
ordre déjà annulé. Une exécution entre les deux intervient pendant que l'annulation est décidée ou en transit, et son rapport
arrive après l'envoi de l'annulation.
\item La fenêtre a pour longueur $L + d$ et les exécutions suivent un processus de Poisson de taux $\lambda$~: la probabilité d'en avoir au moins une vaut
$1 - e^{-\lambda (L + d)}$.
\item $1 - e^{-0.006} \approx 0.60~\%$, $1 - e^{-0.024} \approx 2.4~\%$ et $1 - e^{-0.204} \approx 18.5~\%$.
\item $1 - e^{-x} \approx x - x^2/2$~: l'erreur relative vaut environ $x/2$, inférieure à 5~\% lorsque $\lambda(L + d) < 0.1$.
\item 0.65~\%, 2.1~\% et 11.7~\%.
\item Les ordres sont maintenus une milliseconde en moyenne, soit à peu près la durée du trajet~: beaucoup sont annulés avant d'avoir
séjourné $L + d$ sur la plateforme, et certains qui auraient été en course ont été exécutés avant que la stratégie ne décide d'annuler.
\item La fenêtre cesse de croître à la durée de détention, si bien que la probabilité tend vers celle d'une exécution pendant cette durée~; la
courbe exacte de la figure s'aplatit (16~\% à \qty{3}{\milli\second}), alors que la formule simple continue de croître.
\item $6 / 245 = 2.4~\%$, contre 2.2~\% pour la probabilité exacte~: environ 5.4 courses attendues, avec un écart-type d'environ
2.3, donc oui.
\item L'ancien ordre est exécuté (100), le remplacement est envoyé parce que l'ancien ordre était compté comme disparu, et lui aussi est exécuté~: 200
lots, deux fois la limite.
\item $n$ fois la quantité de l'ordre, soit $100\,n$ lots au-dessus de la limite.
\item Elle continue de compter l'ancien ordre en entier jusqu'à ce que son annulation ou son exécution soit signalée, et refuse le remplacement
d'ici là~: la société ne dépasse jamais la limite, et la stratégie attend un aller-retour avant de pouvoir réutiliser la quantité.
\item $40\,000 \times 2.4~\% \approx 950$ (environ 880 avec la probabilité exacte pour des ordres maintenus une milliseconde).
\item \textbf{Résultat nommé.} Une annulation est en course avec une exécution, avec la probabilité $1 - e^{-\lambda (L + d)} \approx \lambda (L + d)$, pour un
ordre resté au carnet plus longtemps que $L + d$ (2.4~\% à \qty{100}{\micro\second}, 200 exécutions par seconde et \qty{20}{\micro\second} de
délai de rapport)~; sans comptabilisation des \omterm{def:ll:order-gateway-and-order-management:inflight}{ordres en vol} (in-flight), chaque course peut ajouter la quantité d'un ordre de plus à la position, au-delà de la limite.
\item La latence à sens unique $L$ et le délai de rapport $d$ (les deux trajets passant par ses propres systèmes et lignes) et la fréquence de ses annulations~;
non le taux d'exécution, qui dépend du marché.
\item La fenêtre vaut $L + d$~: diviser $L$ par deux ne réduit qu'un seul de ses termes.
\item Les deux parties ont suivi le protocole~: la plateforme a exécuté un ordre actif, puis a répondu correctement que l'annulation était arrivée trop
tard.
\item Des rapports le concernant peuvent encore arriver, et une exécution pour un ordre inconnu serait perdue.
\item Un remplacement entre en course de la même façon, et tant qu'il est en attente, l'ancienne ou la nouvelle quantité peut être exécutée~: la passerelle compte
la plus grande.
\item Compter, à partir du journal, les exécutions arrivées alors que leur ordre était en attente d'annulation, diviser par le nombre d'annulations, et comparer
avec $\lambda (L + d)$ en utilisant les latences mesurées et le taux d'exécution du jour.
\item Compter chaque ordre pour le pire qu'il peut encore faire tant que la plateforme ne dit pas le contraire.
\end{enumerate}
\end{solution}

\begin{solution}{iq:ll:order-gateway-and-order-management:1}
En attente de nouvel ordre, actif, partiellement exécuté, en attente d'annulation, en attente de remplacement, puis exécuté, annulé ou rejeté. La bourse connaît la
vérité~; la firme sait ce qu'elle a demandé et ce qu'elle a reçu en retour, si bien qu'elle a besoin d'états pour «~demandé, pas encore répondu~», pendant lesquels
l'ordre peut encore être exécuté.
\iqlookfor{les états en attente, et pourquoi ils existent.}
\end{solution}

\begin{solution}{iq:ll:order-gateway-and-order-management:2}
L'exécution a eu lieu sur la plateforme avant l'arrivée de l'annulation~: une \omterm{def:ll:order-gateway-and-order-management:inflight}{course annulation--exécution}. Appliquer l'exécution (l'ordre est exécuté ou
partiellement exécuté), s'attendre à ce que l'annulation soit rejetée car trop tardive ou ne porte que sur ce qui reste, et ne jamais traiter l'exécution comme une
erreur, ni l'ordre comme disparu avant le rapport final.
\iqlookfor{la course expliquée, et la comptabilisation dans le pire cas.}
\end{solution}

\begin{solution}{iq:ll:order-gateway-and-order-management:3}
Un \omterm{def:ll:order-gateway-and-order-management:bucket}{seau à jetons} (token bucket) en unités entières, réglé sous la limite de la plateforme, vérifié avant chaque message. Les annulations réduisent le risque, si bien qu'elles ont droit à
une part réservée du budget ou à une priorité sur les nouveaux ordres~; une annulation refusée doit être relancée aussitôt et signalée à la
stratégie.
\iqlookfor{un seau, et des annulations traitées différemment.}
\end{solution}

\begin{solution}{iq:ll:order-gateway-and-order-management:4}
Rien n'est certain~: des ordres ont pu être exécutés ou annulés, par l'annulation à la déconnexion (cancel on disconnect) si elle est activée. Les compter tous en entier, arrêter
les nouveaux ordres, se reconnecter en demandant les rapports manqués, appliquer le rejeu, réconcilier avec le \omterm{def:ll:order-gateway-and-order-management:dropcopy}{drop copy}, et seulement ensuite reprendre.
\iqlookfor{pire cas à l'aveugle, rejeu par numéro de séquence, réconciliation avant la reprise.}
\end{solution}

\begin{solution}{iq:ll:order-gateway-and-order-management:5}
Garder des sommes cumulées de la quantité restante la plus défavorable par côté (et par instrument), mises à jour à chaque modification d'un ordre~: retirer sa
ancienne contribution, appliquer la modification, ajouter la nouvelle. Le contrôle se réduit à une addition et une comparaison.
\iqlookfor{agrégats incrémentaux plutôt que balayages.}
\end{solution}

\begin{solution}{iq:ll:order-gateway-and-order-management:6}
Une passerelle par session de plateforme, avec son automate d'état, son limiteur et son exposition dans le pire cas~; une agrégation à l'échelle de la firme
des expositions au-dessus d'elles pour les limites entre plateformes~; les rapports de la plateforme comme source de vérité pour les exécutions, contrôlés en continu face aux \omterm{def:ll:order-gateway-and-order-management:dropcopy}{drop
copy} et chaque jour face au courtier~; toute rupture est examinée, jamais écrasée.
\iqlookfor{vérité par plateforme issue des rapports, copies indépendantes, et réconciliation.}
\end{solution}
